\chapter{Les matières plastiques}

\noindent\textit{Sacs, bouteilles, jouets, emballages : les matières plastiques sont
partout dans la vie quotidienne. Issues du pétrole, elles posent aujourd'hui un défi
environnemental majeur, ce qui rend indispensable de comprendre leur fabrication et leur
devenir.}
\vspace{8pt}

\connecteBanner

\section{Origine des matières plastiques}

\begin{definition}[Matière plastique]
Une \textbf{matière plastique} est un matériau constitué de \textbf{polymères},
c'est-à-dire de longues molécules formées par l'\textbf{enchaînement répété} d'un grand
nombre de petites molécules identiques appelées \textbf{monomères}. La plupart des
plastiques sont fabriqués à partir de dérivés du pétrole (pétrochimie).
\end{definition}

\begin{propriete}[Polymérisation]
La \textbf{polymérisation} est la réaction chimique au cours de laquelle des monomères
s'assemblent pour former un polymère :
\[
n\ \mathrm{CH_2\!=\!CH_2}\ \longrightarrow\ \mathrm{(-CH_2-CH_2-)}_n
\]
(exemple : l'éthylène, monomère, forme le polyéthylène, polymère, où $n$ désigne un très
grand nombre de répétitions).
\end{propriete}

\begin{illus}
\begin{tikzpicture}[scale=0.85]
\foreach \x in {0,1,2,3,4} {
  \filldraw[onbo!25,draw=onbo,line width=1pt] (\x,1.5) circle(0.35);
}
\node at (2,2.3) {\small monomères};
\draw[->,onbgreen,line width=1.3pt] (2,1) -- (2,0.4);
\node[onbgreen,right] at (2.1,0.7) {\small polymérisation};
\draw[onbmuted,line width=1.5pt] (0,-0.6) -- (0.6,-0.2) -- (1.2,-0.6) -- (1.8,-0.2) --
  (2.4,-0.6) -- (3,-0.2) -- (3.6,-0.6) -- (4.2,-0.2);
\node[below] at (2.1,-0.9) {\small polymère (longue chaîne)};
\end{tikzpicture}
\fig{Figure — Polymérisation : assemblage répété de monomères en une longue chaîne.}
\end{illus}

\section{Les grandes familles de plastiques}

\begin{propriete}[Thermoplastiques]
Les \textbf{thermoplastiques} \textbf{ramollissent sous l'effet de la chaleur} et
peuvent être refondus et remoulés plusieurs fois : ils sont donc \textbf{recyclables}.
Exemples : polyéthylène (PE, sacs, bouteilles), polychlorure de vinyle (PVC, tuyaux),
polytéréphtalate d'éthylène (PET, bouteilles d'eau).
\end{propriete}

\begin{propriete}[Thermodurcissables]
Les \textbf{thermodurcissables} \textbf{durcissent définitivement} lors de leur première
mise en forme (sous l'effet de la chaleur ou d'un durcisseur) : ils ne peuvent
\textbf{plus être refondus}. Exemples : bakélite (prises électriques), résines époxy.
\end{propriete}

\begin{remarque}
Cette différence de comportement à la chaleur explique pourquoi le recyclage des
thermodurcissables est beaucoup plus difficile que celui des thermoplastiques.
\end{remarque}

\section{Propriétés et usages}

\begin{propriete}[Propriétés générales des plastiques]
Les matières plastiques sont généralement \textbf{légères}, \textbf{isolantes}
(électriques et thermiques), \textbf{résistantes à la corrosion} et peu coûteuses à
produire, ce qui explique leur usage massif. Elles présentent aussi des
\textbf{inconvénients} : elles sont pour la plupart \textbf{non biodégradables} et se
dégradent très lentement dans l'environnement (plusieurs centaines d'années).
\end{propriete}

\section{Recyclage et protection de l'environnement}

\begin{definition}[Recyclage]
Le \textbf{recyclage} consiste à collecter, trier, puis retraiter les déchets plastiques
pour en fabriquer de nouveaux objets, plutôt que de les jeter. Un \textbf{symbole
triangulaire numéroté} (de 1 à 7) apposé sur les objets plastiques indique le type de
polymère et facilite le tri.
\end{definition}

\begin{propriete}[Impact environnemental]
Les déchets plastiques mal gérés s'accumulent dans les décharges, les sols et les
océans. Sous l'effet du soleil et des vagues, ils se fragmentent en
\textbf{microplastiques}, ingérés par la faune et présents jusque dans la chaîne
alimentaire humaine.
\end{propriete}

\begin{propriete}[Pistes de réduction]
Réduire l'impact des plastiques passe par : la \textbf{réduction} de leur usage (limiter
les emballages, éviter le plastique à usage unique), le \textbf{tri} et le
\textbf{recyclage} systématiques, et la recherche de \textbf{matériaux alternatifs}
biodégradables.
\end{propriete}

% ═══════════════════════════════════════════════════════════════
\section{L'essentiel à retenir}

\begin{synthese}
\textbf{Polymère.} Longue molécule formée par l'enchaînement répété de monomères
(polymérisation), le plus souvent issus du pétrole.\\[4pt]
\textbf{Thermoplastique.} Ramollit à la chaleur, refondable, recyclable (PE, PVC,
PET).\\[4pt]
\textbf{Thermodurcissable.} Durcit définitivement, non refondable (bakélite, résine
époxy).\\[4pt]
\textbf{Environnement.} Non biodégradables : fragmentation en microplastiques ;
nécessité du tri et du recyclage.\\[4pt]
\textbf{Piège fréquent.} Ne pas confondre « recyclable » (peut être retraité) et
« biodégradable » (se décompose naturellement) : la plupart des plastiques sont
recyclables mais non biodégradables.
\end{synthese}

% ═══════════════════════════════════════════════════════════════
\section{Méthodes \& savoir-faire}

\methode{Distinguer monomère et polymère}
Le \textbf{monomère} est la petite molécule de départ ; le \textbf{polymère} est la
longue chaîne obtenue par répétition de ce monomère.
\begin{exemple}
L'éthylène $\mathrm{CH_2\!=\!CH_2}$ est le monomère du polyéthylène
$\mathrm{(-CH_2-CH_2-)}_n$.
\end{exemple}

\methode{Distinguer thermoplastique et thermodurcissable}
Se demander si le matériau peut être \textbf{refondu} après sa première mise en forme :
oui $\to$ thermoplastique ; non $\to$ thermodurcissable.
\begin{exemple}
Une bouteille en PET peut être fondue et remoulée : c'est un thermoplastique.
\end{exemple}

\methode{Analyser l'impact environnemental d'un objet plastique}
Identifier s'il est à usage unique, s'il est recyclable (symbole triangulaire), et sa
durée de dégradation estimée dans la nature.
\begin{exemple}
Un sac plastique à usage unique, non recyclé, met plusieurs centaines d'années à se
fragmenter dans l'environnement.
\end{exemple}

% ═══════════════════════════════════════════════════════════════
\section{Exercices d'application}
\begin{multicols}{2}

\rubrique{Origine et polymérisation}
\exo{}\diff{1} Qu'est-ce qu'un polymère ?

\exo{}\diff{1} Qu'est-ce qu'un monomère ?

\exo{}\diff{2} À partir de quelle ressource naturelle la plupart des plastiques sont-ils
fabriqués ?

\exo{}\diff{2} Qu'appelle-t-on polymérisation ?

\rubrique{Familles de plastiques}
\exo{}\diff{2} Qu'est-ce qui différencie un thermoplastique d'un thermodurcissable ?

\exo{}\diff{1} Citer deux exemples de thermoplastiques.

\exo{}\diff{1} Citer un exemple de thermodurcissable.

\exo{}\diff{2} Pourquoi les thermodurcissables sont-ils plus difficiles à recycler que
les thermoplastiques ?

\rubrique{Environnement et recyclage}
\exo{}\diff{2} Qu'est-ce qu'un microplastique ?

\exo{}\diff{2} À quoi sert le symbole triangulaire numéroté présent sur les objets
plastiques ?

\exo{}\diff{3} Expliquer la différence entre « recyclable » et « biodégradable ».

% ═══════════════════════════════════════════════════════════════
\end{multicols}

\section{Exercices d'entraînement}
\begin{multicols}{2}

\rubrique{Origine et polymérisation}
\exo{}\diff{2} Écrire l'équation de polymérisation de l'éthylène en polyéthylène.

\exo{}\diff{3} Pourquoi dit-on qu'un polymère est une « très grande molécule » ?

\rubrique{Familles de plastiques}
\exo{}\diff{2} Une pièce plastique fond et peut être remoulée : à quelle famille
appartient-elle ?

\exo{}\diff{2} Une prise électrique en bakélite ne fond pas, même fortement chauffée : à
quelle famille appartient-elle ?

\exo{}\diff{3} Pourquoi les prises électriques sont-elles souvent fabriquées en
matériau thermodurcissable plutôt que thermoplastique ?

\rubrique{Propriétés et usages}
\exo{}\diff{2} Citer deux propriétés qui expliquent l'usage massif des plastiques dans
l'industrie.

\exo{}\diff{3} Pourquoi les plastiques sont-ils souvent utilisés pour isoler les fils
électriques ?

\rubrique{Environnement et recyclage}
\exo{}\diff{2} Citer deux conséquences négatives de l'accumulation de déchets plastiques
dans les océans.

\exo{}\diff{3} Proposer deux gestes concrets permettant de réduire sa consommation de
plastique à usage unique.

\exo{}\diff{3} Expliquer pourquoi le tri des déchets plastiques est une étape essentielle
avant le recyclage.

% ═══════════════════════════════════════════════════════════════
\end{multicols}

\section{Sujets type BEPC}

\sujetbac[Polymères et familles de plastiques]
\begin{enumerate}[label=\arabic*)]
\item Définis les termes « monomère » et « polymère ».
\item Cite un exemple de thermoplastique et un exemple de thermodurcissable.
\item Explique la différence de comportement à la chaleur entre ces deux familles.
\item Laquelle des deux familles est la plus facilement recyclable ? Pourquoi ?
\end{enumerate}

\sujetbac[Plastiques et environnement]
\begin{enumerate}[label=\arabic*)]
\item Pourquoi dit-on que la plupart des matières plastiques ne sont pas
biodégradables ?
\item Qu'est-ce qu'un microplastique et comment se forme-t-il ?
\item Cite une conséquence de la présence de microplastiques dans les océans.
\item Propose deux mesures pour réduire la pollution plastique.
\end{enumerate}

\sujetbac[Propriétés et usages du plastique]
\begin{enumerate}[label=\arabic*)]
\item Cite trois propriétés des matières plastiques qui expliquent leur usage courant.
\item Pourquoi les plastiques sont-ils de bons isolants électriques ?
\item Que signifie le symbole triangulaire numéroté présent sur un emballage plastique ?
\item Quelle est la différence entre recycler et réduire sa consommation de plastique ?
\end{enumerate}

% ═══════════════════════════════════════════════════════════════
\section{Corrigés}
\begin{multicols}{2}

\corrige{de l'exercice 1}
Une longue molécule formée par l'enchaînement répété de petites molécules identiques.

\corrige{de l'exercice 2}
La petite molécule de départ qui se répète pour former un polymère.

\corrige{de l'exercice 3}
Le pétrole (via la pétrochimie).

\corrige{de l'exercice 4}
La réaction chimique au cours de laquelle des monomères s'assemblent pour former un
polymère.

\corrige{de l'exercice 5}
Le thermoplastique ramollit à la chaleur et peut être refondu ; le thermodurcissable
durcit définitivement et ne peut plus être refondu.

\corrige{de l'exercice 6}
Par exemple : le polyéthylène (PE) et le PVC.

\corrige{de l'exercice 7}
Par exemple : la bakélite.

\corrige{de l'exercice 8}
Car ils ne peuvent pas être refondus : leur structure durcie définitivement empêche de
les remouler.

\corrige{de l'exercice 9}
Un petit fragment de plastique (souvent moins de $5$ mm) issu de la fragmentation d'un
déchet plastique dans l'environnement.

\corrige{de l'exercice 10}
Il indique le type de polymère dont est fait l'objet, ce qui facilite son tri en vue du
recyclage.

\corrige{de l'exercice 11}
« Recyclable » signifie que le matériau peut être retraité pour fabriquer un nouvel
objet ; « biodégradable » signifie qu'il se décompose naturellement dans
l'environnement. Un plastique peut être recyclable sans être biodégradable.

\corrige{du sujet 1}
1) Monomère : petite molécule de départ. Polymère : longue molécule formée par
l'enchaînement répété de monomères.\\
2) Par exemple : polyéthylène (thermoplastique), bakélite (thermodurcissable).\\
3) Le thermoplastique ramollit et peut être refondu à la chaleur ; le thermodurcissable
durcit définitivement dès sa première mise en forme.\\
4) Le thermoplastique, car il peut être refondu et remoulé plusieurs fois.

\corrige{du sujet 2}
1) Car leur structure chimique se dégrade extrêmement lentement dans la nature (plusieurs
centaines d'années).\\
2) Un petit fragment de plastique issu de la fragmentation, sous l'effet du soleil et des
vagues, d'un déchet plastique.\\
3) Il peut être ingéré par la faune marine et se retrouver dans la chaîne alimentaire.\\
4) Par exemple : réduire l'usage du plastique à usage unique et développer le
recyclage.

\corrige{du sujet 3}
1) Par exemple : légèreté, résistance à la corrosion, faible coût de production.\\
2) Car ils ne laissent pas passer le courant électrique (matériaux isolants).\\
3) Il indique le type de polymère utilisé et facilite le tri sélectif des déchets.\\
4) Recycler consiste à retraiter un déchet déjà produit ; réduire consiste à limiter la
quantité de plastique consommée à la source.

\end{multicols}
\exercicesBanner
